(a) Suppose $X$ is reduced and $s_x^n=0$. After shrinking the domain of $s$ around $x$, we have $s^n=0$, hence $s=0$ and $s_x=0$. Conversely, if every stalk is reduced and $s^n=0$ on an open set, every germ of $s$ is zero, so $s=0$.

(b) Nilpotents are preserved by restriction, so the indicated quotients form a presheaf. Its associated sheaf has stalk $(\mathcal O_{X,x})_{\mathrm{red}}$. On an affine open $\operatorname{Spec}A$, the quotient map $A\to A_{\mathrm{red}}$ induces a homeomorphism $\operatorname{Spec}A_{\mathrm{red}}\to\operatorname{Spec}A$ and a morphism from this associated sheaf to $\mathcal O_{\operatorname{Spec}A_{\mathrm{red}}}$. On stalks it is the canonical isomorphism $(A_{\mathfrak p})_{\mathrm{red}}\cong(A_{\mathrm{red}})_{\mathfrak p/\sqrt{(0)}}$, so it is an isomorphism. Thus $X_{\mathrm{red}}$ is a scheme. The quotient morphism $\mathcal O_X\to(\mathcal O_X)_{\mathrm{red}}$ is local on stalks and defines the required morphism with underlying map the identity.

(c) Since $X$ is reduced, $f^\#$ kills every nilpotent section. It therefore factors uniquely through the presheaf of reduced rings and hence through its associated sheaf. The resulting stalk maps $(\mathcal O_{Y,f(x)})_{\mathrm{red}}\to\mathcal O_{X,x}$ are local because the original stalk maps are local. This gives the unique required factorization.
